
\subsection{Experiments}
This chapter is about the technical destails of the experiments.
Therefor I will describe the experimental setup and analyse the the arcitecture of the methods, hence the code.
The goal is to supply an indepth understanding of the implementation and the technical details to fulfill a reproduction of my results. 


\subsubsection*{Experimental design}
The experiments are run on a Desktop PC, equiped with a NVIDIA GeForce RTX 3090 and a 12th gen Intel(R) Core(TM) i9 - 12900k 3.2 GHz CPU.
The PC has 65536 RAM and 24326 MB of VRAM. 
It runs Windows 11 Enterprise Version 25H2.
For the code to implement I used the latest version of the integrated development environment (IDE) PyCharm 2025.3.3.
I used python 3.12, because there werent all the packages available for the latest version of python.
I used a virtual environment to manage the dependencies of the project and created a requiremtents.txt file to list all the packages and their versions.


\subsubsection*{InterCont}
The authers of the InterCont method, supplied a reproduction package with their paper. 
The reproduction of the InterCont method went flawlessly and I could reproduce the results of the paper.
Nevertheless, the authers did not include code for other methods, supposably because they are widely used and implemented in standard libraries.
Therefore I implemented the k-NN and MCD methods by myself.


The reproduction package contained 4 different python-files: main, train, data\_loader and helper\_functions.
Beginning with main.py, this file orchestrates the experiment. 
It handles the configuration, load the data and triggers the training process.
Those configurations can be changed by the command line.
I run the code without any command line arguments, so I used the default configurations.
The different configurations are batch size, dataset and faster version.
The default batch size is 3000, which will perform for my relatively small dataset a full batch gradient descent.
The default configuration of 3000 was supposedly chosen by the authers, because in the dataset repository are datasets with more dimensions and instances.
Later I will perform a sensitivity analysis on the batch size, because it is an important hyperparameter for the training process and it can have a significant impact on the results. 
\todo{insert batch size into theory}


The dataset can be changed in the command line argument.
The reproduction package contains all datasets odf the ODDS repository. 
Because I am analysing just one dataset and this dataset happens to be the default setting, I neither use it or change the setting.
The last configuration is the faster version, which I neither use, nor change, because my dataset is relatively small and the default setting is "no" for not implementing the fast version.
 

At the end of the file, the authors implemented a command to ensure that the GPU is primarly used and if it is not avaiable, the processed is been run on the CPU.
Which is also implemented in the train.py file. 
That file determines the training process of the InterCont method.
In the beginning of the file, a datapipeline is been implemented by a datawrapper. 
There the data is determined how many rows a dataset has and how to grab a row. 
For analysing the calculated anomaly score, the wrapper keeps the index names.
The machine learning algorithm then can assign the anomaly score to the right row.


The next encoder class describes the architecture of the method and determines the neural network.
The F network is defined in the following lines beginning with line 28 until line 35.
For the input layer it takes the d-kernel\_size features and puts then through the Tanh activation functionm, resulting in a width of -1 and 1.
The kernel\_size is a hyperparameter, that defined in line 86 until 94. 
Depending on the number of dimensions of the dataset, the kernel\_size is set to 2, 10 or d-150. 
My dataset has 13 dimensions, so the kernel\_size is set to 2.
The hdn\_size refers to the hiddensize argument.
Hiddensize determines how many neurons are in the hidden layers.
Here the hiddenlayers are hardcoded to 200 in line 87. 


The first hidden layer is defined in line 31.
It is defined as an fully connected layer and takes all 200 output values from the previous layer as inputs.
The inputs are projected to double the number of outputs, hence 400 neurons.
Temporarily expanding the dimensionality of the dataset to find non-linear relationships.
The LeakyReLu activation function with its hyperparameter set to 0,2 is implemented in line 32.


The output layer the results from the previous layer is compressed to the predefined hiddensize, hence 200 neurons. 
Because of the compression, the neural network is forced to learn the most important relationsships.
The LeakyReLu is once again implemented.


The smaller G-Network is desigened the same.
From line 38 to line 45 the architecture of the G-Network is defined.
The difference between F and G are in the sizes.
The G network uses the complementary kernel\_size and the number of neurons are divided by 4. 
This network uses only LeakyReLu activation functions.
The hyperparameter stays the same for all activation functions.
The number of neurons are again doubled in the first hidden layer, hence the hiddensize is divided by 2.
\todo{fill in the normalization}

Going ahead, the contrastive 


This method does not implement an cross validation procedure to avoid overfitting. 
Instead they use a permutation based approach to determine the anomaly score.
The dataset is spilt into 50\% training and 50\% testing data.
The training data include only the data from the normal class, while in the test data they include the abnormal data with the normal data. 
This approach is called semi-supervised learning.
To reach the goal of stable and reliable results, the authors implemented a shuffle procedure.
The train and test data stays the same all the time.
But inside this datasets, the order of the columns are being changed, hence shuffeld.

To get randomness into the shuffle process, the authors used the pytorch function torch.randomperm()
This function uses a pseudo-random number generator, hence there is no argument set, the default engine has been used.

The authors tun the code 500 random splits and so did I.
The authors did not provide a loop, so I coded it by myselve.
I did not want to loop the main.py, because then the PC hab to load the data for every loop. 
I applied the semi-supvervised approach in dataloading procedure in the logic of my loop.
I sat the random\_state to 42.
That is a often chosen value in scientific research to ensure reproducability.
I could not tell, if the authors hab chosen another value.  
The metrics are calculated accourdingly to the helper functions.


In the helperfunction\_py is the code for calculating the F1-score.
Here the threshold is not hardcoded as in my other methods.
Instead it is a dynamically coded that meany that the ground thruth lables are incorporated into the threshold calculation.
Then the model takes the threshold and applies it to the sorted outlier ratios.
By comparing the taken instances to the ground truth, the method calculates Precision.
But the number of items picked are equal to the number of actual outliers, Precision equals Recall, hence the F1-score equals the same. 


The authors took k-NN as an classic machine learning algorithm benchmark.
Unfortunately they did not provide in their reproduction package any code to test their results.
Therefore I decided to code my own basic k-NN algorithm to benchmark a machine learning method.
To make the methods comparable I decided to implement a semi-supervised learning approach into the k-NN as well.


\todo{insert arguments why semi supervised can be allpied}

\subsubsection*{k-NN}
The k-NN algorithm is implemented by using the PyOD library, wich is a widely used library for anomaly detection in sientific research.
I also used parts of the well known library scikit-learn to implement the cross validation procedure and to calculate my desired metrics.
First I seperated the dataset into the test and training dataset.
I split the data into same-sized training and testing datasets.
Similiar to the InterCont Method, I put in the normal data into the training dataset and the normal plus abnormal data into the test dataset.


Next the cross validation procedure is implemented by using the KFold function from scikit-learn.
The KFold function takes the number of folds as an argument, which I set to 5.
I chose a 5 fold cross validation approach, because my dataset has not many instances. 
By choosing a relatively small amount of folds, I can ensure that the training dataset is not too small and the testing dataset is not too small as well.
I also set the shuffle argument to True, which means that the data will be shuffled before splitting it into folds.
The random state argument is set to 42, which is a commonly used value in the scientific community to ensure reproducibility of the results.


I sat the contamination parameter to 7,7\%, reflecting the outlier ratio.
That ensures, that the threshold for determining the anomaly score is set to reasonable value grounded on scientif literatur.
The  method parameter determines the distance metric used to calcilate the distances between the data points.
I chose the default setting "largest" which means the distance is calcultate the Minkowski distance with p=2 wich is mathematically equivalent to the Euclidean distance.
The test data is normalized by using the min-max-normalization.


\subsubsection*{MCD}
The FAST-MCD algorithm is also implemented by using the PyOD library. 
I chose an unsupervised learning approach for the MCD algorithm, because it is not possible to implement a semi-supervised learning approach for this method.
Because of its statistical nature, the MCD algorithm does not have a training phase, but it calculates the anomaly score based on the whole dataset.
Therefore I did not implement a cross validation procedure for the MCD algorithm, because it is not possible to split the dataset into a training and testing dataset.
The contamination parameter is also set to the outlier ratio of 7,7\%.
The random state argument is set to 42. 
This ensures that the results are reproducible, because the FAST-MCD algorithm takes random initial subsets and performs the C-Steps thereafter.
By choosing a random state, the reproduction of the results is ensured, because the same random subsets will be chosen in each run of the algorithm.


The authors code is written to implement a high number of different datasets.
My dataset is a .mat-file wich is the dataloader is in part about.
But there 
A lot of code is not necessary to implement my dataset, but I want to keep it, because redundant code will not negatively influence my results.
The interesting thing is, that the hyperparameters are hardcoded and those values are for datasets with more dimensions and instances.
Therefore in the nextchapter I will perform sensitivity analysis on chosen hyperparameters.
Moreover the authors implemented just the F1 and AUC metrics, which I also want to analyse, but I want to analyse the AUPRC as well.
So I made small adjustements to their code which is for transparency reasons marked.






